\section{Protocol}
\label{sec:protocol}

This section walks through TARDUS at the protocol level. The order matches the lifecycle of a coin: setup, deposit, issuance, refresh, withdraw. The fee-payer privacy stack is broken out at the end because it cuts across all spend-class operations.

\subsection{Setup}
\label{sec:proto-setup}

The committee runs a Komlo--Goldberg DKG and ends up holding shares of a joint mint key. The joint public key $pk_{\mathrm{joint}}$, together with the denomination $D$ that it signs and an epoch counter $\mathit{epoch}$, forms a keyset record. The user calls \texttt{RegisterKeyset} on the on-chain program to record this record in the on-chain keyset registry; the call requires at least $t$ committee transcript signatures to be considered valid. The reference implementation supports many keysets per program instance, one per denomination. Common denominations in our devnet deployment are 0.001, 0.005, and 0.01\,SOL.

\subsection{Deposit}
\label{sec:proto-deposit}

A user who wants $D$ lamports' worth of bearer-coin liquidity transfers exactly $D$ lamports to the per-denomination vault PDA (seeds \texttt{["tardus", "vault", D.to\_le\_bytes()]}). The transfer is a vanilla System program transaction; the on-chain TARDUS program is not invoked. This is the only point in the protocol where the user's funding wallet is publicly linked to a TARDUS deposit. After this step, the user has no on-chain link to any coin they subsequently mint.

\subsection{Issuance}
\label{sec:proto-issue}

Issuance is the threshold blind Schnorr ceremony. The user picks a fresh coin secret $x \xleftarrow{\$} \mathbb{Z}_\ell$, derives the candidate public commitment $C_p = x \cdot G$, and blinds it according to the Pointcheval--Stern protocol. Each validator returns a partial blind signature; the user unblinds them, aggregates to a single signature $\sigma$ under $pk_{\mathrm{joint}}$, and verifies that $(C_p, \sigma)$ is valid. The pair $(x, C_p, \sigma, D)$ is the coin. The user stores it locally in an encrypted wallet database.

Issuance happens entirely off-chain. The on-chain surface sees nothing.

\subsection{Refresh}
\label{sec:proto-refresh}

To spend partial value or simply refresh a coin to break linkability, the user runs the cut-and-choose protocol of~\cite{dold2019}. They commit to $\kappa + 1$ candidate coin pairs, each summing to the input coin's denomination. Validators jointly sample a challenge index $r \in [0, \kappa]$, the user opens every pair except the $r$-th, validators verify the opened pairs sum correctly, and only the $r$-th pair is then blind-signed by the committee. The user submits a single on-chain \texttt{Refresh} transaction containing:

\begin{itemize}[leftmargin=2em,itemsep=2pt]
  \item an ed25519 precompile instruction verifying $\sigma$ on the surrendered coin's $C_p$ against the keyset's $pk_{\mathrm{joint}}$;
  \item a \texttt{Refresh} instruction that reads the registry PDA, computes $\mathit{nullifier} = \operatorname{SHA-256}(\mathtt{"TARDUS-nullifier-v1"} \,\|\, C_p)$, inserts it into the on-chain nullifier set, and returns success.
\end{itemize}

The new coins are bearer instruments held off-chain. The on-chain effect of a Refresh is exactly one nullifier insertion. Observers see neither the new coins nor any link to the surrendered one.

\subsection{Withdraw}
\label{sec:proto-withdraw}

A Withdraw transaction looks like a Refresh but moves real SOL out of the vault. The user supplies a coin $(C_p, \sigma, D)$ and a recipient pubkey. The program verifies the signature via the same ed25519 precompile, computes the nullifier, checks it is not already in the set, inserts it, and uses \texttt{invoke\_signed} to transfer $D$ lamports from the vault PDA to the recipient. The recipient pubkey is pinned in the instruction data, so the transaction signer cannot redirect funds en route.

We use Withdraw to close the economic loop on devnet: a fresh recipient wallet receives native SOL through a TARDUS coin that was minted via the off-chain blind-sign ceremony, delivered via a sealed-box payload through our relay, and surrendered on chain to release the vault funds. The fee payer of the Withdraw need not be the recipient, and (per the fee-payer privacy stack below) typically is not.

\subsection{Sealed-box delivery}
\label{sec:proto-sealedbox}

To hand a coin from Alice to Bob without the on-chain surface seeing the hand-off, we use a small relay daemon that stores opaque ciphertexts keyed by recipient pubkey. Alice computes a sealed-box AEAD ciphertext for Bob's receiving pubkey and POSTs it to the relay; Bob polls for messages addressed to his pubkey, fetches the ciphertext, and decrypts.

The construction mirrors NaCl's \texttt{crypto\_box\_seal}~\cite{nacl} but ends with an Ed25519 receiving pubkey rather than a native X25519 one. The wire format is:
\[
  \mathit{wire} = \underbrace{epk^{(M)}}_{\text{32 bytes}} \,\|\, \underbrace{\operatorname{ChaCha20Poly1305}_{K, N}(\mathit{plaintext})}_{\text{$|m| + 16$ bytes}},
\]
where $(K, N) = \operatorname{HKDF\text{-}SHA\text{-}256}(\mathit{salt} = \mathtt{"TARDUS-sealed-box-v1"}, \mathit{ikm} = esk \cdot pk_R^{(M)}, \mathit{info} = epk^{(M)} \,\|\, pk_R^{(M)})$ and the ephemeral keypair $(esk, epk^{(M)})$ is freshly sampled by the sender.

The receiver, holding the Ed25519 secret scalar $sk$, recovers the same shared secret by computing $sk \cdot epk^{(M)}$ on the Montgomery curve. Notice this requires the scalar multiplication to be unclamped: the BIP-39-derived $sk$ is a canonical Ed25519 scalar modulo $\ell$, not an X25519-clamped value per RFC 7748. The distinction matters for any implementation, because most off-the-shelf X25519 libraries clamp by default and will not produce a shared secret that decrypts a TARDUS sealed-box. We state this explicitly in the specification appendix so independent implementations do not rediscover the constraint by failing round-trips.

\subsection{Fee-payer privacy stack}
\label{sec:proto-feepayer}

Solana's base layer reveals every transaction's \texttt{fee\_payer} pubkey. If the user always pays their own fees, a graph crawler can correlate every spend they ever make. The fee-payer stack neutralises this in four composable layers.

\textbf{Layer 1 --- ephemeral payer.} For each spend-class transaction, the client generates a fresh ed25519 keypair, transfers a small amount of SOL into it from a sponsor wallet (default $10^{-3}$\,SOL), and uses the ephemeral as the transaction's \texttt{fee\_payer}. The ephemeral never signs anything else and is dropped after the transaction. Each TARDUS spend then has a previously-unseen signer.

\textbf{Layer 2 --- multi-sponsor rotation.} The single sponsor of Layer 1 is replaced by a small set of sponsor keypairs, picked at random per ephemeral. This breaks the "all ephemerals funded from the same wallet" residual link, at the cost of operating multiple sponsor wallets.

\textbf{Layer 3 --- on-chain sponsor pool.} A program-owned PDA accepts deposits from anyone and pays out fixed amounts (default $10^{-3}$\,SOL) to fresh ephemeral keypairs. The pool commingles all deposits, so the on-chain link from an ephemeral payer back to "the wallet that funded my fees" runs through the pool's aggregated balance rather than through a direct transfer. The pool is rate-limited at one payout per five Solana slots ($\approx 2.5$\,s) to bound drain attacks at roughly 14\,SOL per hour under default parameters.

\textbf{Layer 4 --- pool caller rotation.} The wallet that submits the \texttt{SponsorPayout} transaction is itself drawn at random from the sponsor pool, closing the "the deployer always calls payout" residual leak.

The composed stack is the subject of Theorem T9 (\Cref{thm:T9}). On devnet we have $|\mathcal{F}| = 5$ sponsor wallets and the direct-link probability sits at roughly $20\%$; on mainnet, with $|\mathcal{F}| \geq 100$ under healthy deposit cadence, the bound drops below $1\%$.
